\chapter{Architettura della libreria}
\label{sec:architecture}

\section{Pipeline}

Una \textit{pipeline}, in informatica, è una sequenza di stadi (\textit{stage}) attraverso cui i dati fluiscono in modo ordinato. Ciascuno \textit{stage} riceve un input, lo elabora e produce un output che diventa l’input dello \textit{stage} successivo. Ogni \textit{stage} ha una responsabilità ben delimitata, favorendo così modularità, riuso e facilità di ragionamento sull’intero sistema.

In contesti concorrenti, ogni \textit{stage} può operare in parallelo sugli elementi che riceve. Ciò significa che, mentre lo \textit{stage} 1 elabora un nuovo messaggio, lo \textit{stage} 2 può già processare quello precedente, incrementando così il \textit{throughput} complessivo della pipeline.

All’interno della libreria \textbf{goccia}, il concetto di \textit{pipeline} è incarnato dalla struct \texttt{Pipeline}, che funge da orchestratore degli \textit{stage} e rappresenta il punto di ingresso per l’utilizzo della libreria stessa. Questa struct definisce i tre metodi principali che ogni \textit{stage} deve implementare: \texttt{Init}, \texttt{Run} e \texttt{Close}, corrispondenti alle diverse fasi del ciclo di vita della pipeline (si veda la sezione successiva per ulteriori dettagli). Inoltre, il metodo \texttt{AddStage} consente di aggiungere nuovi \textit{stage} fintanto che la pipeline non è in esecuzione.

<disegno pipeline>

\section{Stage}

Uno \textit{stage} di una pipeline rappresenta un’unità modulare di elaborazione che riceve un input da \textit{stage} precedenti, oppure da una sorgente esterna, lo elabora e produce un output per gli \textit{stage} successivi. In tal modo, gli \textit{stage} formano una catena di trasformazione dei dati scalabile e facilmente componibile. Dal punto di vista concettuale, uno \textit{stage} costituisce un limite logico e sequenziale all’interno della \textit{pipeline}, con un flusso dati che attraversa le fasi di input, elaborazione e output, in maniera simile a quanto avviene nei processi CI/CD (build, test, deploy).

Nella libreria \textbf{goccia}, gli \textit{stage} devono aderire all’interfaccia \texttt{Stage}, la quale definisce i metodi necessari per rendere una struct di Go compatibile con la pipeline principale. Tali metodi coincidono con le tre fasi del ciclo di vita: \texttt{Init}, per l’inizializzazione delle risorse; \texttt{Run}, per l’elaborazione dei dati in input; e \texttt{Close}, per il rilascio delle risorse e la chiusura dello \textit{stage}.

La libreria categorizza gli \textit{stage} in tre gruppi principali — \texttt{Ingress}, \texttt{Processor} ed \texttt{Egress} — tutti aderenti all’interfaccia generica \texttt{Stage}. Un aspetto rilevante riguarda la modalità di esecuzione: il metodo \texttt{Run} viene invocato all’interno di una goroutine dedicata per ciascuno \textit{stage}. Di conseguenza, una pipeline composta da cinque \textit{stage} disporrà di almeno cinque goroutine attive, oltre a quella principale di esecuzione nel main.

<disegno lifecycle>

\section{Connector}

Affinché una \textit{pipeline} possa effettivamente operare, una volta creati gli \textit{stage}, è necessario metterli in comunicazione tramite un componente capace di trasferire i dati da uno \textit{stage} al successivo. Tale componente deve essere \textit{thread safe}, poiché ogni \textit{stage} viene eseguito su una goroutine distinta, e deve disporre di un meccanismo di buffer per mitigare gli effetti della \textit{backpressure}, fenomeno che si manifesta quando uno \textit{stage} a valle elabora i dati più lentamente del precedente, causando un rallentamento a catena.

Il linguaggio Go fornisce, a tal fine, i channel (anche bufferizzati) come tipo primitivo per la comunicazione tra goroutine. Tuttavia, la libreria \textbf{goccia} adotta un approccio differente, basato su un \textit{ring buffer} lock-free, con l’obiettivo di massimizzare il \textit{throughput}. In particolare, si è optato per un’implementazione \textit{Single Producer Single Consumer} (SPSC), poiché la relazione tra gli \textit{stage} della pipeline è di tipo 1:1. La versione \textit{Multiple Producer Multiple Consumer} (MPMC) del \textit{ring buffer} viene invece impiegata nel contesto del \textit{worker pool} (descritto successivamente).

<tabella channel vs ring buffer>

Per mantenere la libreria estensibile a futuri sviluppi, viene definita l’interfaccia \texttt{Connector}, dotata di quattro metodi principali:
\begin{itemize}
    \item \texttt{Write}: per scrivere un dato nel connettore;
    \item \texttt{Read}: per leggere il dato disponibile;
    \item \texttt{Close}: per chiudere il connettore e impedire ulteriori scritture;
    \item \texttt{SetReadTimeout}: per definire un tempo massimo di attesa in caso di assenza di dati in ingresso.
\end{itemize}

\section{Worker Pool}

Un \textit{worker pool} rappresenta una strategia per gestire l’esecuzione concorrente di numerosi compiti (\textit{task}) mediante un numero limitato di lavoratori (\textit{worker}) riutilizzabili nel tempo. Questo approccio consente di evitare la creazione eccessiva di unità di lavoro parallele, situazione che potrebbe portare alla saturazione delle risorse di sistema.

In genere, un \textit{worker pool} è composto da un numero fisso di \textit{worker}, determinato dal programmatore o configurato a \textit{runtime} tramite una variabile d’ambiente. Nel linguaggio Go, la funzione \texttt{runtime.NumCPU()} permette di ottenere il numero di core fisici e virtuali disponibili, favorendo così una gestione ottimizzata delle risorse. Questo rappresenta il punto di partenza per implementare un \textit{worker pool} efficiente, indipendentemente dall’hardware sottostante.

La libreria \textbf{goccia} estende tale concetto introducendo uno \textit{scaler} automatico in grado di aumentare o ridurre dinamicamente il numero di \textit{worker} in base alla quantità di \textit{task} presenti nella coda da cui essi estraggono il lavoro. Tale coda effettua un’operazione di \textit{fan-out}, distribuendo i task in ingresso ai diversi \textit{worker}, e funge anche da buffer intermedio. A valle dei \textit{worker}, un componente analogo svolge la funzione di \textit{fan-in}, ovvero la ricomposizione sequenziale dei risultati derivanti dal processamento dei \textit{task}. Entrambi i componenti sono implementati mediante \textit{ring buffer} MPMC.

Lo \textit{scaler} valuta periodicamente il numero di \textit{worker} attivi, basandosi sulla dimensione della coda di \textit{fan-out}. Se il numero di \textit{task} supera la soglia di riferimento (\texttt{QueueDepthPerWorker}), viene segnalato al \textit{worker pool} di aggiungere un nuovo \textit{worker}; altrimenti, ne viene rimosso uno. Qualora lo \textit{scaler} rilevi ripetutamente la necessità di ridurre i \textit{worker}, viene applicato un tempo di \textit{backoff} esponenziale tra le riduzioni successive. In questo modo, il \textit{worker pool} reagisce rapidamente a un improvviso incremento di carico, riducendo poi progressivamente le risorse una volta gestito il picco.

<disegno worker pool con fan-in/fan-out>


\section{Ingress Stage}

Uno \textit{Ingress stage} rappresenta il punto di ingresso dei dati all'interno della \textit{pipeline} di elaborazione. Costituisce la prima componente che riceve gli input provenienti dall'esterno del sistema e li introduce nel flusso di elaborazione, alimentando così l'intera catena di \textit{stage} successivi. In sostanza, uno \textit{Ingress stage} agisce come una sorgente di dati, traducendo eventi, messaggi o pacchetti provenienti da differenti protocolli o interfacce in un formato gestibile dalla \textit{pipeline} di \textbf{goccia}.

Nel contesto della libreria, gli \textit{Ingress stage} sono progettati per essere altamente performanti, concorrenti e facilmente integrabili con i meccanismi di comunicazione già definiti dal framework. Ogni \textit{Ingress stage} è quindi responsabile di stabilire una connessione con la sorgente, acquisire i dati in ingresso, trasformarli nel formato previsto e inoltrarli verso lo \textit{stage} successivo.

La libreria \textbf{goccia} fornisce diversi tipi di \textit{Ingress stage}, ciascuno progettato per scenari operativi o fonti di dati differenti.

\subsection{Ticker}

Lo \textit{Ticker stage} rappresenta l'\textit{Ingress stage} più semplice messo a disposizione dalla libreria \textbf{goccia} ed è concepito per generare in modo autonomo un flusso regolare di messaggi, senza dipendere da sorgenti esterne quali socket di rete, file o code di messaggistica. Dal punto di vista concettuale, si tratta di una sorgente periodica di eventi, particolarmente utile per casi d'uso quali il testing della \textit{pipeline}, il benchmarking di prestazioni, l'iniezione di traffico sintetico o la generazione di "heartbeat" applicativi a intervalli costanti. Ogni "tick" corrisponde alla produzione di un messaggio che attraversa l'intera \textit{pipeline}, permettendo di osservare il comportamento dei vari \textit{stage} a valle in presenza di un carico controllato.

La logica di generazione dei messaggi sfrutta il \texttt{time.Ticker} della libreria standard del Go, il quale mette a disposizione un channel notificato ogni volta che l'intervallo $T$ viene raggiunto, dove $T$ è definito dalla configurazione dello \textit{stage}.

\begin{lstlisting}[language=Go, caption={Ciclo principale del Ticker stage (ingress/ticker.go, righe 90-103)}]
for {
	// ...
	select {
	case <-ctx.Done():
		return
	case <-ts.ticker.C:
		// generate the ticker message
	}
}
\end{lstlisting}

Nel frammento di codice riportato, è possibile osservare come il costrutto \texttt{select} permetta di ascoltare su molteplici channel in modalità multiplex. In questo caso specifico, il \textit{runtime} rimane in attesa di essere notificato dal ticker oppure della ricezione di un segnale di cancellazione del contesto (\texttt{ctx.Done()}). Quest'ultimo meccanismo è fondamentale per implementare strategie intelligenti di rilascio delle risorse, ad esempio quando il processo riceve un segnale \texttt{SIGINT} o \texttt{SIGTERM}, consentendo così l'implementazione di uno \textit{graceful shutdown} coerente nelle applicazioni. Per questa ragione, nelle librerie Go moderne è prassi richiedere un \texttt{context.Context} come primo argomento di funzioni che impiegano socket o richiedono tempistiche significative per la terminazione.

Il tipo di messaggio prodotto dallo \textit{Ticker stage} è descritto dalla struct \texttt{TickerMessage} e contiene un unico campo specifico, \texttt{TickNumber}, utilizzato per numerare progressivamente i tick. Al momento della generazione del messaggio vengono popolati i campi generici di metadato (tempo di ricezione, timestamp) e viene inizializzata la \textit{root span} per la tracciatura degli attraversamenti degli \textit{stage} successivi, conformemente allo standard OpenTelemetry\cite{otel:trace:api}.

\subsection{UDP}

Lo \textit{UDP stage} rappresenta un \textit{Ingress stage} progettato per ricevere datagrammi \textit{UDP} provenienti dalla rete. A differenza dello \textit{Ticker stage}, che genera autonomamente messaggi a intervalli regolari, lo \textit{UDP stage} rimane in ascolto su un socket di rete e resta in attesa di nuovi datagrammi in arrivo. Questa tipologia di \textit{Ingress stage} risulta particolarmente utile in contesti dove è necessario elaborare flussi di dati provenienti da sorgenti esterne, quali sensori IoT, applicazioni remote, strumenti di monitoraggio di rete o qualsiasi sistema che comunichi tramite il protocollo \textit{UDP}. La natura \textit{connectionless} del protocollo \textit{UDP} consente di ricevere messaggi da molteplici mittenti senza la necessità di stabilire connessioni esplicite, rendendolo ideale per scenari ad alto \textit{throughput} dove la perdita occasionale di datagrammi è accettabile in cambio della bassa latenza.

La libreria \textbf{goccia} mette a disposizione tre parametri di configurazione per lo \textit{UDP stage}: indirizzo IP, porta e dimensione del buffer. I primi due servono a impostare la sorgente da cui ricevere i pacchetti e dispongono di valori di default — \texttt{"0.0.0.0"} per l'indirizzo (che indica l'ascolto su tutte le interfacce di rete) e \texttt{20.000} per la porta. La dimensione del buffer utilizzato per leggere il payload dei datagrammi \textit{UDP} è impostata di default a $1474$ byte, poiché la dimensione massima per un payload Ethernet standard è di $1500$ byte, da cui si sottraggono $28$ byte relativi all'header \textit{UDP}.

Internamente, lo \textit{stage} utilizza un puntatore a una connessione \textit{UDP} fornita dalla libreria standard del Go\cite{go:net:udpconn}. Tale connessione viene inizializzata nel metodo \texttt{Init} del ciclo di vita dello \textit{stage}. Una volta che l'inizializzazione ha esito positivo, è possibile procedere con l'esecuzione dello \textit{stage} tramite il metodo \texttt{Run}. Questo metodo implementa un ciclo \texttt{for} in cui viene richiamato il metodo \texttt{Read} (operazione bloccante) della connessione \textit{UDP}, il quale scrive i byte ricevuti in un buffer pre-allocato alla dimensione configurata. Parallelamente al ciclo principale, viene avviata una goroutine ausiliaria il cui compito è chiudere la connessione \textit{UDP} nel momento in cui viene ricevuta una notifica di cancellazione tramite \texttt{context.Context}. Una volta che la connessione è chiusa, il metodo \texttt{Read} ritorna un errore di tipo \texttt{net.ErrClosed}\cite{go:net:errclosed}, permettendo così al ciclo di uscire in modo ordinato.

\begin{lstlisting}[language=Go, caption={Ciclo principale dello UDP stage (udp.go, metodo run)}]
go func() {
	<-ctx.Done()
	us.conn.Close()
}()

buf := make([]byte, us.bufferSize)

for {
	// Read the UDP payload
	n, err := us.conn.Read(buf)
	if err != nil {
		// Check if the connection is closed
		if errors.Is(err, net.ErrClosed) {
			// Check if caused by context cancellation
			select {
			case <-ctx.Done():
				return
			default:
			}
		}

		// ...
	}

	// Handle the buffer and send the message
}
\end{lstlisting}

Il tipo di messaggio prodotto dallo \textit{UDP stage} è descritto dalla struct \texttt{UDPMessage}, contenente due campi specifici: \texttt{Payload} e \texttt{PayloadSize}. Il messaggio implementa l'interfaccia \texttt{message.Serializable} definita nel package \texttt{message}, consentendo il collegamento di differenti tipi di \textit{Processor stage} a valle che accettano messaggi recanti buffer di byte. Al fine di soddisfare l'interfaccia, è sufficiente implementare il metodo \texttt{GetBytes}, il quale restituisce una slice di byte corrispondente al campo \texttt{Payload} del messaggio.

\begin{lstlisting}[language=Go, caption={Metodo GetBytes di UDPMessage (ingress/udp.go)}]
func (um *UDPMessage) GetBytes() []byte {
	return um.Payload
}
\end{lstlisting}

Poiché il payload è caratterizzato da dimensioni significative, al fine di ridurre il carico sul \textit{garbage collector} attraverso il riuso della memoria e l'eliminazione di continue allocazioni e deallocazioni di migliaia di oggetti, si è optato per l'utilizzo di un \textit{object pool} (\texttt{sync.Pool})\cite{go:sync:pool}. Gli oggetti \texttt{UDPMessage} vengono quindi pre-allocati nel pool e riutilizzati per ogni nuovo datagramma ricevuto, riducendo così la pressione sul sistema di gestione della memoria durante l'elaborazione ad alto throughput.

Lo \textit{stage} espone due metriche destinate al monitoraggio del numero di messaggi e byte ricevuti, implementate tramite l'uso di contatori asincroni\cite{otel:metrics:asynccounter}. Tali metriche possono essere impiegate per monitorare il \textit{throughput} in ingresso alla \textit{pipeline}, fornendo visibilità sulla velocità di ricezione dei dati dalla sorgente di rete.


\begin{itemize}
    \item TCP
    \item Kafka
    \item File
    \item eBPF
\end{itemize}

\section[Stage di processamento]{Stage di processamento}

\begin{itemize}
    \item Cannelloni encoder/decoder
    \item CAN (acmelib)
    \item CSV encoder/decoder
    \item Custom
    \item Filter
    \item Tee
    \item ROB
\end{itemize}

\section[Stage di egress]{Stage di egress}

\begin{itemize}
    \item UDP
    \item TCP
    \item Kafka
    \item File
    \item QuestDB
    \item Sink
\end{itemize}

\section[Observabilità con OpenTelemetry]{Observabilità con OpenTelemetry}

\begin{itemize}
    \item Metriche
    \item Distributed tracing
\end{itemize}